\documentclass[11pt]{article}
\usepackage[a4paper,margin=27mm]{geometry}
\usepackage{amsmath,amssymb,mathtools}
\usepackage{fontspec,unicode-math}
\setmainfont{Latin Modern Roman}
\setmathfont{Latin Modern Math}
\usepackage{xurl}
\usepackage[hidelinks]{hyperref}
\newcommand{\sref}[2]{\hyperlink{source:#1:#2}{[S#2]}}
\newcommand{\cataloguescope}[1]{\paragraph{Mathematical question.}#1}
\hypersetup{pdftitle={UnsolvedMath Problem 9400124}}
\begin{document}
% UnsolvedMath ID 9400124; source catalogue code AMR-093-0124
\setcounter{section}{9400123}
\section{Fortune's conjecture}\label{9400124}
{\small\sffamily UnsolvedMath ID 9400124 · Number Theory}
\subsection{Definitions and mathematical statement}

\paragraph{Shared notation.}$\mathbb N=\{1,2,\ldots\}$, and $\mathbb Z,\mathbb Q,\mathbb R,\mathbb C$ denote the integers, rationals, reals and complex numbers. Any other conventions are specified locally.


Let $p_1=2<p_2=3<\cdots$ be the primes and let $P_n=\prod_{j=1}^{n}p_j$ for $n\ge1$. Define the $n$th Fortunate number by
\[
F_n=\min\{m\in\mathbb N:m>1,\ P_n+m\text{ is prime}\}.
\]
The set is nonempty because primes are unbounded. The condition $m>1$ is part of the definition.
\paragraph{Statement recorded in the supplied source.}
Is $F_n$ prime for every $n\ge1$? \sref{9400124}{2}
\subsection{Short English statement}
Is the smallest increment greater than one which turns each primorial into a prime always itself prime?
\subsection{Sources}
\begin{description}
\item[\hypertarget{source:9400124:1}{S1}] ULAM.AI and the UnsolvedMath Contributors, UnsolvedMath: A Curated Collection of Open Mathematics Problems, record 9400124 (AMR-093-0124). CC BY 4.0. \url{https://huggingface.co/datasets/ulamai/UnsolvedMath}.
\item[\hypertarget{source:9400124:2}{S2}] Wikipedia contributors, \emph{List of unsolved problems in mathematics}, revision 1366636767 (2026) Number theory, Fortune's conjecture. \url{https://en.wikipedia.org/w/index.php?title=List_of_unsolved_problems_in_mathematics&oldid=1366636767#Number_theory}.
\item[\hypertarget{source:9400124:3}{S3}] OEIS Foundation and contributors, \emph{The On-Line Encyclopedia of Integer Sequences}, A005235: Fortunate numbers: least m > 1 such that m + prime(n)\# is prime, where p\# denotes the product of all primes <= p., official data snapshot \texttt{693b5499c9d2} (2026) Name, definition and comments. \url{https://raw.githubusercontent.com/oeis/oeisdata/693b5499c9d248f570dd2d979e8bd1d7eecf689b/seq/A005/A005235.seq}.
\end{description}
\label{end:9400124}
\end{document}
